% from ega4-18-fr.tex
\oldpage[IV-4]{196}

En vertu de (19.2.4), la condition $a)$ équivaut à dire que $B$ est un $A$-module plat et que $\mathfrak J\otimes_A k$ (qui s'identifie à un idéal de $C\otimes_A k$ en vertu de la platitude de $B$ sur $A$ ($0_{\mathrm I}$, 6.1.2)) est régulier dans cet anneau. Comme l'anneau $C\otimes_A k$ est régulier, il revient au même, en vertu de (19.3.2), de dire (lorsque $B$ est un $A$-module plat) que $\mathfrak J\otimes_A k$ est un idéal régulier de $C\otimes_A k$, ou que $B\otimes_A k=(C\otimes_A k)/(\mathfrak J\otimes_A k)$ est un anneau d'intersection complète; d'où le corollaire.

\begin{definition}[19.3.6]
\EGAReaderTarget{ega4.c04.def.19-3-6}{IV.19.3.6}
Soit $f:X\to S$ un morphisme plat localement de présentation finie. On dit que $X$ est une \emph{intersection complète relativement à $S$ en un point} $x$, si la fibre $f^{-1}(f(x))$ est une intersection complète (absolue) au point $x$. On dit que $X$ est une \emph{intersection complète relativement à $S$}, ou que le morphisme $f$ est un \emph{morphisme d'intersection complète}, si $X$ est intersection complète relativement à $S$ en chacun de ses points.
\end{definition}

\begin{proposition}[19.3.7]
\label{IV.19.3.7}
Soient $g:X\to S$, $h:Y\to S$ deux morphismes plats localement de présentation finie, $f:Y\to X$ une $S$-immersion, $y$ un point de $Y$, $x=f(y)$, $s=g(x)=h(y)$. On suppose que la fibre $X_s=g^{-1}(s)$ soit régulière au point $x$. Alors, pour que $f$ soit une immersion transversalement régulière au point $y$, il faut et il suffit que $Y$ soit une intersection complète relativement à $S$ au point $y$.
\end{proposition}

Dire que $f$ est transversalement régulière au point $y$ revient, en vertu de (19.2.4), à dire que $f_s:Y_s\to X_s$ est une immersion régulière au point $y$. Mais comme $X_s$ est régulier au point $x=f(y)$, cela équivaut, en vertu de (19.3.2), à dire que $Y_s$ est intersection complète (absolue) au point $y$, d'où la proposition.

\begin{corollary}[19.3.8]
\label{IV.19.3.8}
Soit $f:X\to S$ un morphisme plat localement de présentation finie. L'ensemble $U$ des $x\in X$ en lesquels $X$ est une intersection complète relativement à $S$ est ouvert dans $X$. Si de plus $f$ est propre, l'ensemble $E$ des $s\in S$ tels que $X_s=f^{-1}(s)$ soit intersection complète en chacun de ses points est une partie ouverte de $S$.
\end{corollary}

Comme $E=Y-f(X-U)$, la seconde assertion résulte de la première et du fait que $f$ est une application fermée. La première assertion est locale sur $X$, donc on peut supposer que $X$ est un sous-préschéma fermé d'un préschéma de la forme $Z=S[T_1,\ldots,T_n]$ ($T_i$ indéterminées). Ce dernier étant lisse sur $S$ (17.3.8), ses fibres $Z_s$ sont régulières, et il revient donc au même, en vertu de (19.3.7), de dire que $x\in U$ ou que l'immersion $X\to Z$ est transversalement régulière au point $x$; on en conclut donc le corollaire par (19.2.2).

\begin{proposition}[19.3.9]
\label{IV.19.3.9}
\begin{enumerate}
  \item[(i)] Toute immersion ouverte est un morphisme d'intersection complète.
  \item[(ii)] Soit $f:X\to S$ un morphisme plat et localement de présentation finie qui est un morphisme d'intersection complète. Pour tout changement de base $g:S'\to S$, $f'=f_{(S')}:X_{(S')}\to S'$ est un morphisme d'intersection complète. La réciproque est vraie si $g$ est fidèlement plat et quasi-compact.
  \item[(iii)] Si $f:X\to Y$ et $g:Y\to Z$ sont deux morphismes plats localement de présentation finie qui sont des morphismes d'intersection complète, il en est de même de $g\circ f:X\to Z$.
\end{enumerate}
\end{proposition}

L'assertion (i) découle aussitôt de la définition (19.3.1). Pour prouver (ii), remarquons d'abord que si $f$ est plat et localement de présentation finie, il en est de

